
\documentclass[11pt]{article}

\usepackage[colorlinks=true,linkcolor=black,citecolor=black,urlcolor=black]{hyperref}
\usepackage{amsmath,mathpazo,graphicx,doi}

\usepackage[a4paper,margin=3cm]{geometry}

\usepackage{xspace}
\usepackage{setspace}


\raggedright
%%\RequirePackage{lineno} 
%%\modulolinenumbers[2]

\usepackage[style=nature,articletitle=true,isbn=false,backend=bibtex]{biblatex}
\bibliography{editorial}
\AtEveryBibitem{\clearfield{month}}
%% Following bibtex hack overcomes  an error that came in recently to
%% biblatex in 2017.
%%http://tex.stackexchange.com/questions/311426/bibliography-error-use-of-blxbblverbaddi-doesnt-match-its-definition-ve
\makeatletter
\def\blx@maxline{77}
\makeatother


\begin{document}

\onehalfspacing

%% PNG taken from https://orcid.org/sites/default/files/images/orcid_128x128.png
\newcommand{\orcid}[1]{\textsuperscript{\href{https://orcid.org/#1}
  {\includegraphics[width=2.5mm]{orcid_128x128.png}}}}

\title{Towards standard practices for sharing computer code and programs in neuroscience}

\author{\orcid{0000-0001-8607-8025}Stephen J. Eglen$^{1,\dagger}$, 
  \orcid{0000-0001-7879-4531}Ben Marwick${^2}$, 
  \orcid{0000-0003-3456-2493}Yaroslav O. Halchenko${^3}$, 
  \orcid{0000-0001-6398-6370}Michael~Hanke${^{4,5}}$,
  \\
  \orcid{0000-0001-6390-2616}Shoaib Sufi$^6$,
  \orcid{0000-0001-5963-8576}Padraig Gleeson$^7$, 
  \orcid{0000-0002-5480-6638}R. Angus Silver$^7$,
  \orcid{0000-0002-4793-7541}Andrew P. Davison$^8$,
  \\
  \orcid{0000-0002-0885-1825}Linda Lanyon$^9$,
  \orcid{0000-0001-9438-9691}Mathew Abrams$^9$,
  \orcid{0000-0003-2015-6590}Thomas Wachtler$^{10}$,
  \\
  \orcid{0000-0003-3814-7529}David J. Willshaw$^{11}$,
  \orcid{0000-0002-2844-8099}Christophe Pouzat$^{12}$,
  %\\
  \orcid{0000-0002-9794-749X}Jean-Baptiste Poline$^{13,\dagger}$}
\date{\today}
\maketitle

\noindent $^{1}$ 
Cambridge Computational Biology Institute,
Department of Applied Mathematics and Theoretical Physics,
University of Cambridge, UK
%%Wilberforce Road, Cambridge CB3 0WA, UK

\noindent $^{2}$ 
Department of Anthropology,
University of Washington,
Seattle, WA 98195-3100 USA 

\noindent $^{3}$
Department of Psychological and Brain Sciences,
Dartmouth College,
Hanover, NH 03755 USA

\noindent $^{4}$
Institute of Psychology II, Otto-von-Guericke-University Magdeburg, 39106 Magdeburg, Germany

\noindent $^{5}$
Center for Behavioral Brain Sciences, 39106 Magdeburg, Germany

\noindent $^{6}$
Software Sustainability Institute, University of Manchester, UK

\noindent $^{7}$
Department of Neuroscience, Physiology and Pharmacology, 
University College London, UK


\noindent $^{8}$
Unit\'{e} de Neurosciences, Information et Complexit\'{e}, CNRS, Gif sur Yvette, France


\noindent $^{9}$
International Neuroinformatics Coordinating Facility, 
Karolinska Institutet, Stockholm, Sweden

\noindent $^{10}$
Department of Biology II, Ludwig-Maximilians-Universit\"{a}t
M\"{u}nchen, Germany

\noindent $^{11}$
Institute for Adaptive and Neural Computation, 
School of Informatics, University of Edinburgh, UK


\noindent $^{12}$ MAP5 Paris-Descartes University and CNRS UMR 8145, 
45 rue des Saints-P\`{e}res, 75006 Paris, France

\noindent $^{13}$  Henry H. Wheeler, Jr. Brain Imaging Center, Helen Wills Neuroscience Institute, University of California, Berkeley, USA

\vspace*{2mm}
\noindent $\dagger$
Corresponding authors: \url{S.J.Eglen@damtp.cam.ac.uk}, \url{jbpoline@gmail.com}

\clearpage

\begin{abstract}
  Computational techniques are central in many areas of neuroscience,
  and are relatively easy to share.  This paper describes why computer
  programs underlying scientific publications should be shared, and
  lists simple steps for sharing.  Together with ongoing efforts in
  data sharing, this should aid reproducibility of research.
\end{abstract}

\renewcommand{\cite}[1]{\autocite{#1}}

%%\linenumbers

\section*{Background}

Many areas of neuroscience are now critically dependent on
computational tools to help understand the large volumes of data being
created.  Furthermore, computer models are increasingly being used to
help predict and understand the function of the nervous system.  Many
of these computations are complex and usually cannot be concisely
reported in the methods section of a scientific article.  In a few
areas there are widely used software packages for
analysis (e.g., SPM, FSL, AFNI, FreeSurfer, Civet in
neuroimaging) or simulation (e.g. NEURON, NEST, Brian).
However, we often write new computer programs to solve
specific problems in the course of our research.  Some of these
programs may be relatively small scripts that help analyze all of our
data, and these rarely get described in papers.  As authors, how best
can we maximize the chances that other scientists can reproduce our
computations, find errors, or reuse our methods on their data?  Is our research
reproducible \cite{web:naturefocus}?

To date, the sharing of computer programs underlying neuroscience
research has been the exception (see below for some examples), rather
than the rule.  However, there are many potential benefits to sharing
these programs, including increased understanding and reuse of your
work.  Furthermore, open source programs can be scrutinized and
improved, whereas the functioning of closed source programs remains
forever unclear \cite{Vihinen2015}.  Funding agencies, research
institutes and publishers are all gradually developing policies to
reduce the withholding of computer programs relating to research
\cite{Morin2012-65e}.  The Nature family of journals has published
opinion pieces in favor of sharing whatever code is available, in
whatever form \cite{Barnes2010-iv,Ince2012-225}.  Since October 2014,
all Nature journals require papers to include a statement declaring
\textit{whether} the programs underlying central results in a paper
are available. In April 2015 \textit{Nature Biotechnology} offered
recommendations for providing code with papers and began asking
referees to give feedback on their ability to test code that
accompanies submitted manuscripts \cite{NatBiotech2015}.  In July 2015
F1000Research stated that ``Software papers describing non-open
software, code and/or web tools will be rejected''
(\url{http://f1000research.com/channels/f1000-faculty-reviews/for-authors/article-guidelines/software-tool-articles}). Also
in July 2015, BioMed Central introduced a minimum standards of
reporting checklist for BMC Neuroscience and several other journals,
requiring submissions to include a code availability statement and for
code to be cited using a DOI or similar unique identifier
\cite{Kenall2015}.  We believe that all journals should adopt policies
that highly encourage, or even mandate, the sharing of software
relating to journal publications as this is the only practical way to
check the validity of the work.


\subsection*{What should be shared?}

It may not be obvious what to share, especially for complex projects
with many collaborators.  As advocated by Claerbout and Donoho, for
computational sciences the scholarship is not the article; the
"scholarship is the complete software [...]"
\cite{claerbout_electronic_1992,donoho_invitation_2010}.  So, ideally,
we should share all code and data needed to allow others to
reproduce our work, but this may not be possible or practical.
However, it is expected that the key parts of the work should be shared,
e.g. implementations of novel algorithms or analyses.  At a minimum,
we suggest following the recommendation of submission of work to
ModelDB\cite{modeldb}, i.e. to share enough code, data and
documentation to allow at least one key figure from your manuscript to
be reproduced.  However, by adopting appropriate software tools, as
mentioned in the next section, it is now relatively straightforward to
share the materials required to regenerate \textit{all} figures and
tables.  Code that already exists, is well tested and documented, and is reused 
in the analysis should be cited. Ideally, all other code should be 
communicated, including code that performs simple
preprocessing or statistical tests, or code that deals with local
computing issues such as hardware and software configurations. 
While this code may not be reusable, it will help others understand how
analyses are performed, find potential mistakes, and aid reproducibility. 
Finally, if the work is computationally intensive and requires a long
time to run (e.g. many weeks), one may prefer to provide a small ``toy''
example to demonstrate the code.

By getting into the habit of sharing as much as possible, not only do
we help others who wish to reproduce our work (which is a basic tenet
of the scientific method), we will be helping other members of our
laboratory, or even ourselves in the future.  By sharing our code
publicly, we are more likely to write higher-quality code
\cite{Easterbrook2014}, and we will know where to find it after we
have moved on from the project \cite{Halchenko2015}, rather than the
code disappearing on a colleague's laptop when they leave your group,
or suffer some misfortune \cite{Berg2016}.  We also will be part of a
community and benefit from the code shared by others, thus reducing
software development time for ourselves and others.

\subsection*{Simple steps to help you share code}

Once you have decided \textit{what} to share, here are some simple
guidelines for \textit{how} to share the work.  Ideally, these
principles should be followed throughout the lifetime of the research
project, not just at the end when we wish to publish our
results. Guidelines similar to these have been proposed in many areas
of science \cite{Nosek2015, miguel2014, stodden2012journals},
suggesting that they are part of norms that are emerging across
disciplines. In the `further reading' section below, we list some
specific proposals from other fields that expand on the guidelines we
suggest here.

\begin{description}
\item [Version control] Use a version control system (such as Git) to
  develop the code \cite{Blischak2016}.  The version control repository
  can then be easily and freely shared with others using sites such as
  \url{http://github.com} \cite{Ram2013} or
  \url{https://bitbucket.org}.  These sites allow you fine control
  over private versus public access to your code.  This means that you
  can keep your code repository private during its development, and
  then publicly share the repository at a later stage e.g. at the time
  of publication, although we recommend opening the code from the
  start of the project.  It also makes it easy for others to
  contribute to your code, and to adapt it for their own uses.

\item [Persistent URLs] Generate stable URLs (such as a DOI) for key
  versions of your software.  Unique identifiers are a key element in
  demonstrating the integrity and reproducibility of research
  \cite{vasilevsky2013reproducibility}, and allow referencing of the
  exact version of your code used to produce figures. DOIs can be
  obtained freely and routinely with sites such as
  \url{http://zenodo.org} and \url{http://figshare.com}.  If your work
  includes computer models of neural systems, you may wish to consider
  depositing these models in established repositories such as
  ModelDB\cite{modeldb}, Open Source Brain\cite{osb} or NITRC
  \cite{poline_software_2014}. Some of these sites allow for private
  sharing of repositories with anonymous peer reviewers. Journal
  articles that include a persistent URL to code deposited in a
  trusted repository meet the requirements of level two of the
  `analytic methods (code) transparency' standard of the TOP
  guidelines \cite{Nosek2015}.


\item [License] Choose a suitable license for your code to assert how
  you wish others to reuse your code.  For example, to maximize reuse,
  you may wish to use a permissive license such as MIT or BSD
  \cite{stodden2009enabling}.  Licenses are also important to protect you from
  others misusing your code.  Visit \url{http://choosealicense.com/}
  to get a simple overview of which license to choose, or 
  \url{http://www.software.ac.uk/resources/guides/adopting-open-source-licence}
  for a detailed guide.

\item [Etiquette] When working with code written by others, observe
  Daniel Kahneman's 'reproducibility etiquette'\cite{Kahneman2014} and
  have a discussion with the authors of the code to give them a chance
  to fix bugs or respond to issues you have identified before you make
  any public statements.  Cite their code in an appropriate fashion.


\item [Documentation] Contrary to popular expectations, you do not
  need to write extensive documentation or a user's guide for the code
  to still be useful to others \cite{Barnes2010-iv}.  However, it is
  worth providing a minimal README file to describe what the code
  does, and how to run it.  For example, you should provide
  instructions on how to regenerate key results, or a particular
  figure from a paper. Literate programming methods, where code and
  narrative text are interwoven in the same document, make
  documentation semi-automatic and can save a lot of time when
  preparing code to accompany a publication \cite{schulte2012multi,
    gentleman2012statistical}. However, these methods admittedly take
  more time to write in the first instance, and you should be prepared
  to rewrite documentation when rewriting code.  In any cases,
  well-documented code allows for easier re-use and checking.

\item [Tools] Consider using modern, widely used software tools that
  can help with making your computational research reproducible.  Many
  of these tools have already been used in neuroscience and serve as
  good examples to follow, for example Org mode \cite{Delescluse2011},
  IPython/Jupyter \cite{Stevens2013} and Knitr \cite{Eglen2014}.
  Virtualization environments, such as VirtualBox appliances and
  Docker containers, can also be used to encapsulate or preserve all
  of the computational environment so that other users can run your
  code without having to install numerous dependencies
  \cite{Boettiger2015}.


\item [Case studies] In addition to the examples listed above in
  Tools\cite{Delescluse2011,Stevens2013,Eglen2014}, there are many
  prior examples to follow when sharing your code.  For example, some
  prominent examples of reproducible research in computational
  neuroscience include Vogels et al. \cite{Vogels2011-c8c} and Waskom
  et al. \cite{Waskom2014-gd}; see
  \url{https://github.com/WagnerLabPapers} for details.  The ModelDB
  repository contains over 1000 computational models deposited with
  instructions for reproducing key figures to papers
  e.g. \url{https://senselab.med.yale.edu/ModelDB/showModel.cshtml?model=93321}
  for a model of activity-dependent conductances \cite{Liu1998}.



\item [Data] Any experimental data collected alongside the software
  should also be released or made available.  For small datasets, this
  could be stored alongside the software, although it may be
  preferable to store experimental data separately in an appropriate
  repository.  Both PLOS and Scientific Data maintain useful lists of
  subject-specific and general repositories for data storage, see
  \url{http://journals.plos.org/plosbiology/s/data-availability#loc-recommended-repositories}
  and \url{http://www.nature.com/sdata/data-policies/repositories}.
  

  
\item [Standards] Use of (community) standards, where appropriate, should
  be encouraged, in particular use of non-proprietary formats
  to enable long-term accessibility.
  In computational neuroscience for example, PyNN
  \cite{Davison2009} and NeuroML \cite{CannonEtAl2014} are widely used
  formats for making models more accessible and portable across
  multiple simulators.  Neuroimaging data and results can be organized
  using BIDS \cite{Gorgolewski2016}.


\item [Tests] Testing the code has long been recognized
  as a critical step in the software industry but the practice is not
  widely adopted yet by researchers. We recommend including test
  suites that demonstrate the code is producing the correct
  results\cite{Axelrod2014-xi}. These tests can be at a low level
  (testing each individual function, called unit testing) or at a
  higher level (e.g. testing that the program yields correct answers
  on simulated data) \cite{wilson_best_2014}. With public data available,
  it is often straightforward to have a test verifying that published
  results can be recomputed. Linking tests to continuous 
  integration services (such as Travis CI, \url{https://travis-ci.org}) allows 
  these tests to be automatically run each time a change is made to the code,
  ensuring failing tests are immediately flagged and can be dealt with quickly.

\item [User support] Although some people are eager to provide support
  for their code after it has been published, others may feel that
  they do not want to be burdened by e.g. feature requests.  One
  simple suggestion to avoid this is to establish a user community for
  the code \cite{Gorgolewski2016b}.  This could be as simple as
  creating a mailing list or asking for issues to be posted on a
  github repository.

\end{description}

\subsection*{Closing remarks}

Changing the behaviors of neuroscientists so that they make their code more
available will likely be resisted by those who do not see the community
benefits as outweighing the personal costs of the time and effort required to
share code \cite{stodden2010scientific}. The community benefits, in our view,
are obvious and substantial: we can demonstrate more robustly and transparently
the reliability of our results, we can more easily adapt methods developed by
others to our data, and the impact of our work increases as others can
similarly reuse our methods on their data. Thus, we will endeavor to lead by
example, and follow all these practices as part of our future work in all
scientific publications. Even if the code we produce today will not run ten
years from now, it will still be a more precise and complete expression of our
analysis than the text of the methods section in our paper. 

However, exhortations such as this article are only a small part of
making code sharing a normal part of doing neuroscience; many other
activities are important.  All researchers should be trained in sound
coding principles; such training is provided by organizations such as
Software Carpentry \cite{wilson_best_2014} or Data Carpentry and
through national neuroinformatics initiatives,
e.g. \url{http://python.g-node.org}.  Furthermore, we should request
code and data when reviewing, and submit to and review for journals
that support code sharing.  Grant proposals should be checked for
mentions of code availability, and we should encourage efforts toward
openness in hiring, promotion, and reference letters
\cite{leveque2012reproducible}. Funding agencies and editors should
also consider mandating code sharing by default.  This combination of
efforts on a variety of fronts will increase the visibility of
research accompanied by open source code, and demonstrate to others in
the discipline that code sharing is a desirable activity that helps
move the field forward.

We believe that the sociological barriers to code sharing are harder
to overcome than the technical ones. Currently, academic success is
strongly linked to publications and there is little recognition for
producing and sharing code. Code may also be seen as providing a
private competitive advantage to researchers. We challenge this view
and propose that code be regarded as part of the research products and
part of the publication in which should be shared by default, and that
there should be an obligation to share code for those conducting
publicly funded research. We hope the code availability review
\cite{nn2017ed} will help establish such sharing as the norm.
Moreover, we are advocating for code sharing as part of a broader
culture change embracing transparency, reproducibility, and
re-usability of research products.

\subsection*{Acknowledgments}

This article is based upon discussions from a workshop to encourage
sharing in neuroscience, held in Cambridge, December 2014.  It was
financially supported and organized by the International
Neuroinformatics Coordinating Facility (\url{http://www.incf.org}),
with additional support from the Software Sustainability institute
(\url{http://www.software.ac.uk}).
%
Michael Hanke was supported by funds from the German federal state of
Saxony-Anhalt and the European Regional Development Fund (ERDF), Project:
Center for Behavioral Brain Sciences.

\subsection*{Author Contributions}
All authors contributed to discussions and writing/editing the manuscript.


\subsection*{Competing Financial Interests Statement}
None.

\printbibliography


\clearpage

\subsection*{Box 1: Further reading}

Khodiyar, V. 2015. Code Sharing --- read our tips and share your own. Scientific Data Blog, February 19, 2015. \url{http://blogs.nature.com/scientificdata/2015/02/19/code-sharing-tips/}

Kitzes, J., Turek, D., \& Deniz, F. (Eds.). 2017. The Practice of Reproducible Research: Case Studies and Lessons from the Data-Intensive Sciences. Oakland, CA: University of California Press. \url{https://www.practicereproducibleresearch.org/}

Leveque, R. 2013. Top ten reasons to not share your code (and why you
should anyway). SIAM News, April 2013, \url{https://sinews.siam.org/Details-Page/top-ten-reasons-to-not-share-your-code-and-why-you-should-anyway}


Stodden, V., M. McNutt, D. H. Bailey, E. Deelman, Y. Gil, B. Hanson, M. A. Heroux, J.P. A. Ioannidis and M. Taufer 2016. Enhancing reproducibility for computational methods. Science 354(6317):1240.  DOI: \url{http://dx.doi.org/10.1126/science.aah6168}

Stodden V., \& Miguez, S., 2014. Best practices for computational science: software infrastructure and environments for reproducible and extensible research. Journal of Open Research Software. 2(1), p.e21. DOI: \url{http://dx.doi.org/10.5334/jors.ay}

Stodden, V., Leisch, F., \& Peng, R. (Eds.). 2014. Implementing reproducible research. CRC press, Chapman and Hall. \url{https://osf.io/s9tya/}

Halchenko, Y. O. and Hanke, M. 2015. Four aspects to make science open
``by design'' and not as an after-thought. GigaScience, 4. DOI:
\url{http://dx.doi.org/10.1186/s13742-015-0072-7}

Sandve, G. K., Nekrutenko,  A., Taylor,  J., \& Hovig E 2013. Ten simple rules for reproducible computational research. PLoS Comput Biol 9:e1003285.

\clearpage

\subsection*{Box 2: Online communities discussing code sharing}

\begin{description}

\item [StackExchange and related projects] StackExchange is a network of free and highly active question-and-answer websites.  Two members of the network are relevant to questions of code sharing: \url{http://stackoverflow.com/} which is dedicated to questions about programming in any language in any context, and \url{http://academia.stackexchange.com/questions/tagged/reproducible-research} which is focused questions relating to reproducible research in academic context. A related project is \url{https://neurostars.org/} which is a similar free public Q\&A website focused on neuroinformatics questions, and with many questions on software packages, etc. 

\item [Scientists for Reproducible Research] This is an international multi-disciplinary email list that discusses a wide range of issues relating to code sharing: \url{https://groups.google.com/forum/\#!forum/reproducible-research}

\item [GitHub] GitHub is an online repository for computer code and programs that has a large community of researchers that develop and share their code openly on the site. GitHub is the largest and most active code sharing site (others include BitBucket and GitLab) and has convenient tools for facilitating efficient collaborative coding \cite{tippmann2014my,Perez-Riverol2016}. If you are using an open source program you may find a community of users and developers active on GitHub, where you can ask questions and report problems.

\end{description}



\end{document}



%%% Local Variables:
%%% mode: latex
%%% TeX-PDF-mode: t
%%% TeX-master: t
%%% ispell-local-dictionary: "american"
%%% auto-fill-inhibit-regexp: ".*[&|].*[&|].*"
%%% End:

%  LocalWords:  reproducibility DOI neuroinformatics GitHub
